% Generated from chapter-05.md - do not edit.

\chapter{Thinking About The Three Principles}%
\index{three principles}\nobreak%

Why are we talking about the Three Principles in a book called What's a Better Thought? Because Thought is enormously important in the way we experience being alive.

What we think about something can affect how we feel about it, what we see in it, and what we do next.

The Three Principles offer a way of understanding just how powerful and ever-present Thought is in our lives. And that helps explain why finding a better thought can sometimes be one of the simplest ways to find a better feeling, a better perspective, and a little more peace and happiness.

\section{A Kinder Way of Understanding How We Experience Life}%
\index{kinder}\nobreak%

Every now and then, someone explains something in a way that feels both simple and profound at the same time. That's how I felt when I first heard about what are called the \textbf{Three Principles} --- originally named \textbf{Mind, Consciousness, and Thought}.\index{principles}

They were introduced by a quiet, thoughtful man named \textbf{Sydney Banks},\index{Sydney Banks} who had no formal background in psychology or academia. His realization was simple and life-changing:\\
\textbf{everything we experience in life comes from the inside out, not from the world itself.}\\
What's happening ``out there'' doesn't directly create our feelings --- our inner process does.

That insight opened up a whole new way of seeing life for me.

Traditionally, Sydney describes the Three Principles like this:

\begin{itemize}
\item \textbf{Mind} --- the universal intelligence behind life\index{mind}
\item \textbf{Consciousness} --- our capacity to be aware
\item \textbf{Thought} --- the creative power that forms our moment-to-moment experience
\end{itemize}

I honor that language and where it came from.

And---for my own understanding, and for the sake of being grounded and practical---I think about it with slightly different terms:

\textbf{The Brain, Consciousness, and Thought are working together to create our entire human experience.}

\section{Brain}%
\index{brain}\index{mind}\nobreak%

\textbf{For me, the brain is where the miracle of human experience shows up in everyday life.}

I am a living human being with a remarkable brain---an organ so powerful and compassionate that it quietly runs my breathing, digestion, heartbeat, immune system, and nervous system . . . 

and creates my thinking, my emotions, my perceptions, my joys, and my worries.

\emph{(All without asking me to fill out any forms.)}

In this book, I sometimes use the words brain and mind somewhat interchangeably when I am talking about thinking, remembering, imagining, and experiencing life. I know they are not technically the same thing: the brain is the physical organ, while the mind refers more broadly to our thoughts and conscious experience. My purpose here is not to settle that distinction, but to explore the very practical ways our thinking affects how we experience our lives.

\section{Consciousness --- The Gift of Awareness}%
\index{consciousness}\index{awareness}\nobreak%

Consciousness is what allows us to be aware at all.

It's the light that brings our thoughts to life.

It's what lets us notice a beautiful sunset, feel comfort from a warm cup of tea, experience fear, or suddenly see something from a new angle.\index{focus}

It's the difference between simply existing . . .  and experiencing.\index{experiencing}

You know how a dream feels completely real while you're inside it?

That's consciousness at work.

\emph{(Even the very strange dreams.)}

\section{Pause for a Moment}%
\index{exercises}\index{practices}\index{action steps}\nobreak%

Right now, what are you aware of?

A sound?\\
A feeling in your body?\\
A thought passing through?

Notice how awareness brings it to life.

And gently consider:

\textbf{``I'm feeling this way because I'm conscious of my thinking---not because life itself is actually terrible.''}

\section{Thought --- The Creative Power of Experience}

Thought is the moment-to-moment creative force of our lives.\index{creative force}

We use it constantly---usually without noticing---to create our inner world.

Every feeling, every story, every judgment, every worry, every hope . . . 

we experience them through the lens of thought.

We don't live in the feeling of our circumstances.

We live in the feeling of our thinking about our circumstances.

\emph{(Which explains why two people can experience the same event . . .  and feel completely different.)}

\textbf{This doesn't mean thoughts are bad or that we should try to control them.}

\textbf{It simply means they are always changing---like weather in the sky.}

\emph{(Sometimes sunny, sometimes stormy, occasionally a bit unpredictable.)}

\textbf{Better thought:}

\emph{``This is just a passing thought---not a permanent truth.''}

\section{Why This Is So Comforting}%
\index{comforting}\nobreak%

What comforts me most in understanding all of this is something very gentle:

\textbf{We are not broken.}\\
\textbf{We don't need fixing.}

We are simply living inside the ever-changing flow of brain, consciousness, and thought.

Sometimes we create storms in our own inner weather.

\textbf{And sometimes we forget that the sky always clears.}

\emph{(Even when it takes a little while.)}

When we realize that our feelings are coming from the thoughts we're having---not directly from the world---it becomes easier to be kind to ourselves.

To pause.
To breathe.
To wait for the storm to pass.

And often, all it takes is one small shift---

\textbf{just one better thought---}

\textbf{to change everything.}

Throughout this book, I emphasize the powerful relationship between thought and emotion. I believe that the meaning we give to what happens---our interpretations, expectations, memories, and thoughts---plays an enormous part in what we feel.

Emotions can sometimes seem to arise before we are aware of having a thought at all. Our brains and bodies make very rapid assessments of what is happening, often outside our conscious awareness. So when I speak of Thought coming before feeling, I am using the word broadly to include these quick interpretations and meanings as well as the conscious thoughts we can actually hear ourselves thinking.

My purpose is not to claim that every emotion begins with a clearly identifiable sentence in the mind. It is to explore something very practical: when we become aware of what we are thinking and find a different way to understand what is happening, our emotional experience can change too. That possibility is at the heart of What's a Better Thought?

\section{A Gentle Noticing About Beliefs}%
\index{belief}\nobreak%

I've also begun to notice something important:

\textbf{Thoughts that are repeated often---and felt strongly---can become beliefs.}

\textbf{They start as passing ideas . . . }

\textbf{but over time, they can feel solid, fixed, and true.}

\textbf{And yet---even beliefs began as thoughts.}

\emph{(Very convincing thoughts . . .  but still thoughts.)}

\textbf{Which means:}

\textbf{They can be questioned.}\\
\textbf{They can soften.}\\
\textbf{They can change.}

Chapter 6 takes a closer look at Beliefs.

\section{A Quiet Ending Thought}

You don't have to understand all of this perfectly.

You don't have to ``get it right.''

Just begin to notice:\index{practices}\index{exercises}\index{action steps}

\begin{itemize}
\item What am I thinking right now?
\item How is that thought shaping how I feel?
\end{itemize}

And if you'd like, gently ask:

\textbf{What's a better thought?}

\emph{Even a small one can begin to shift your whole experience.}
